% !TEX root = ../main.tex
% =====================================================================
% 正文第 1 章
%
% 以下文字与附件 3 正文第 1 章完全一致，逐字照录（含附件 3 自带的格式说明
% 段落）。章标题、节标题、小节标题的分级依附件 3 所标示的字号：
%   黑体小二号（18pt）→ 章标题
%   宋体四号（14pt）  → 节标题（\section）
%   宋体小四（12pt）  → 小节标题（\subsection）
%
% 章标题为附件 3 原文“1  **”，故用不编号章命令原样输出，
% 以保证目录条目与正文标题和附件 3 逐字一致。
% =====================================================================

\ujschapter{1\quad 绪论}

\section{课题背景和意义}

\subsection{研究背景}

说明本课题的意义、目的、研究范围及要达到的技术要求；简述本课题在国内外的发展概况及存在的问题；说明本课题的指导思想；阐述本课题应解决的主要问题。

\subsection{研究意义}

……………………（正文小四，中文宋体，英文数字Times New Roman，1.5行距，首行缩进2字符）

\section{课题研究现状}

\section{本文研究内容}

正文部分：章标题为黑体小二号字体，居中打印，段前、段后间距为1行。

二级标题为宋体四号字，加粗，左对齐，段前、段后间距为0.5行。

三级标题为宋体小四号字，加粗，左对齐，段前、段后间距为0.5行。

正文中汉字用宋体小四号，英文和数字用Times New Roman小四号，1.5倍行距，首行缩进2字符。每一个新的章节均另起一页开始书写。

% !TEX root = ../main.tex
% 格式测试用引用，文献来自官方模板示例，与论文主题无关。
\section{参考文献引用格式测试}
以下段落仅用于检查 BibTeX 引用和参考文献排版，正式写作时可删除本节及对应的 input。

单篇引用测试：中文专著的引用应显示为上标方括号编号\upcite{chen2000}。再次引用同一文献\upcite{chen2000}，编号应保持一致。

连续多篇引用测试：一次引用三篇文献\upcite{harrison2012,yu2019,zhongguo1985}，检查连续编号是否合并为范围。

非连续引用测试：引用已经出现过的第一篇与第三篇文献\upcite{yu2019,chen2000}，检查编号是否按升序排列，并用逗号分隔。

平排引用测试：讨论文献\cite{chen2000}时，编号应与正文处于同一基线；句末上标引用应位于标点之前\upcite{chen2000}。

文献类型测试：英文会议论文集\upcite{babu2014}，中文报纸\upcite{lidalu1998}，英文报纸\upcite{french1987}，中文期刊\upcite{chen1996}，英文期刊\upcite{caplan1993}。

学位论文和报告测试：中文学位论文\upcite{zhang2020}，英文学位论文\upcite{calms1965}，中文报告\upcite{feng1997}，机构作者报告\upcite{who1970}，在线报告\upcite{guowuyuan2013}。

其他类型测试：标准\upcite{quanguo2010}，专利\upcite{liujialin1993}，专著析出文献\upcite{cheng1999}，会议析出文献\upcite{jia2011}，电子公告\upcite{guowuyuan2014,xiaoniu2001,dublincore2012}。参考文献表应按首次引用顺序列出全部 21 条文献，并检查长标题和 URL 的换行。

